\documentclass{report}
\usepackage{import}
\usepackage{tcolorbox}
\usepackage{amsmath}
\usepackage{amssymb}
\usepackage{amsthm}
\usepackage{geometry}
\usepackage{enumitem}
\usepackage{pdfpages}
\usepackage{transparent}
\usepackage[hidelinks]{hyperref}
\usepackage{booktabs}
\usepackage{bm}
\usepackage{tabularx}
\tcbuselibrary{theorems}
\tcbuselibrary{skins}

% tabularx Settings
\renewcommand\tabularxcolumn[1]{>{\centering\arraybackslash}m{#1}}

% Page Settings
\newgeometry{vmargin={15mm}, hmargin={15mm, 15mm}}
\setlength{\parindent}{0pt}
\setlength{\parskip}{1em plus 1pt minus 1pt}

% hyperref Settings
\hypersetup{
    colorlinks=true,
    linkcolor=blue,
    filecolor=magenta,      
    urlcolor=cyan,
}

% inkscape Settings
\newcommand{\incfig}[1]{%
    \def\svgwidth{\columnwidth}
    \import{./figure/}{#1.pdf_tex}
}

% Definition Box Styles
\newtcbtheorem[auto counter, number within=subsection]{defn}{Definition}{
% basic settings
    theorem name and number,
    separator sign={},
    description delimiters={(}{)},
    enhanced jigsaw,
    sharp corners,
% box margins and rules
    boxrule=1pt,
    left=0.2cm,right=0.2cm,top=0.2cm,
    toptitle=0.1cm,
    bottomtitle=0.08cm,
% box colors
    colframe=white!25!black,
    colback=white,
    coltitle=white,
% box fonts
    fonttitle=\bfseries,fontupper=\normalsize
}{defn}

% Theorem Box Styles
\newtcbtheorem[auto counter, number within=subsection]{thm}{Theorem}{
% basic settings
    theorem name and number,
    separator sign={},
    description delimiters={(}{)},
    enhanced jigsaw,
    sharp corners,
% box margins and rules
    boxrule=1pt,
    left=0.2cm,right=0.2cm,top=0.2cm,
    toptitle=0.1cm,
    bottomtitle=0.08cm,
% box colors
    colframe=white!25!black,
    colback=white,
    coltitle=white,
% box fonts
    fonttitle=\bfseries,fontupper=\normalsize
}{thm}

% Lemma Box Styles
\newtcbtheorem[number within=tcb@cnt@thm]{lem}{Lemma}{
% basic settings
    theorem name and number,
    separator sign={},
    description delimiters={(}{)},
    enhanced jigsaw,
    sharp corners,
% box margins and rules
    boxrule=1pt,
    left=0.2cm,right=0.2cm,top=0.2cm,
    toptitle=0.1cm,
    bottomtitle=0.08cm,
% box colors
    colframe=white!50!black,
    colback=white,
    coltitle=white,
% box fonts
    fonttitle=\bfseries,fontupper=\normalsize
}{lem}

% Corollary Box Styles
\newtcbtheorem[number within=tcb@cnt@thm]{cor}{Corollary}{
% basic settings    
    theorem name and number,
    separator sign={},
    description delimiters={(}{)},
    enhanced jigsaw,
    sharp corners,
% box margins and rules
    boxrule=1pt,
    left=0.2cm,right=0.2cm,top=0.2cm,
    toptitle=0.1cm,
    bottomtitle=0.08cm,
% box colors
    colframe=white!50!black,
    colback=white,
    coltitle=white,
% box fonts
    fonttitle=\bfseries,fontupper=\normalsize
}{cor}

% Remark Box Styles
\newtcbtheorem[number within=subsection]{rem}{Remark}{
% basic settings
    theorem name and number,
    separator sign={},
    description delimiters={(}{)},
    enhanced jigsaw,
    sharp corners,
% box colors
    colframe=white!60!black,
    colback=white,
    coltitle=black,
% box margins and rules
    boxrule=1pt,
    left=0.2cm,right=0.2cm,top=0.2cm,
    toptitle=0.1cm,
    bottomtitle=0.08cm,
% box fonts
    fonttitle=\bfseries,fontupper=\normalsize
}{rem}

% Proof Box Styles

\title{Preliminaries of Machine Learning}
\author{Hyoung Joon, Kim}
\date{\today}


\begin{document}

\maketitle
\tableofcontents
\newpage

\chapter{Introduction}

This is some notes about preliminaries of machine learning, based on the book by Kevin Murphy.

    \section{Notation}
    $\mathbb{E}_{x \sim p(x)}$ means that the expectation is taken about $x$ where $x$ has distribution $p(x)$.

\chapter{Probability}
\section{Basic Definitions}


% definition of probability space
\begin{defn}{Probability Space}{PSpace}
    A \emph{probability space} is a triple \(( \Omega, \mathcal{F}, \mathbb{P})\) where
    \begin{itemize}
        \item \(\Omega\) is a non-empty set called the \emph{sample space};
        \item \(\mathcal{F}\) is a \(\sigma\)-algebra of subsets of the sample space \(\Omega\);
        \item \(\mathbb{P} :\mathcal{F} \to [0, 1]\) is a probability measure.
    \end{itemize}
\end{defn}

\section{Lists of Probability Distributions}

\subsection{Discrete Distributions}

\begin{center}
%\begin{tabular}
\begin{tabularx}{\textwidth}{c X c c c}
    \toprule
    \textbf{Name} & \textbf{Support} & \textbf{Parameters} & \textbf{PMF} & \textbf{Mean, Variance} \\
    \midrule
    Bernoulli & \(\{0, 1\}\) & \(p \in \{0, 1\}\) & Ber(\(x|p\)) = \(p^x (1-p)^{1-x}\) & \(\mathbb{E}[X] = p, \mathbb{V}[X] = p(1-p)\) \\
    \midrule
    Binomial & \(\{0, 1, 2, \ldots, N\}\) & \(N, p \in [0,1]\) & Bin(\(x|N,p\)) = \(\binom{N}{x} p^x (1-p)^{N-x}\) & \(\mathbb{E}[X] = Np, \mathbb{V}[X] = Np(1-p)\) \\
    \midrule
    Categorical & \(\{0, 1, 2, \ldots, N\}\) & \(\param\) & Cat(\(x|\param\)) = \(\prod_{k=1}^{K}\theta_{k}^{\mathbb{I}(x=k)}\) \\
    \midrule
    Multinomial & \(\bm{x}\) where \(\sum_{k = 1}^{K} x_k = N\) & \(N, \param\) & Mult(\(\bm{x}|\param\)) = \(\binom{N}{x_1, x_2, \ldots, x_K}\prod_{k = 1}^{K}\theta_{k}^{x_k}\) \\
    \midrule
    Poisson & \(\{0, 1, 2, \ldots\}\) & \(\lambda > 0\) & Pois(\(x|\lambda\)) = \(e^{-\lambda}\dfrac{\lambda^x}{x!}\) & \(\mathbb{E}[X] = \lambda, \mathbb{V}[X] = \lambda\) \\
    \bottomrule
\end{tabularx}
%\end{tabular}
\end{center}


\subsection{Continuous Distributions}

\begin{center}
\begin{tabularx}{\textwidth}{c c c X X}
    \toprule
    \textbf{Name} & \textbf{Support} & \textbf{Parameters} & \textbf{PDF} & \textbf{Mean, Variance} \\
    \midrule
    Gaussian &
    \(\mathbb R\) &
    \(\mu, \sigma\) &
    \(\mathcal N(\mu, \sigma^2)=\dfrac{1}{\sqrt{2\pi\sigma}}e^{-\frac{1}{2\sigma^2}(x-\mu)^2}\) &
    \(\mathbb{E}[X] = \mu, \mathbb{V}[X] = \sigma^2\)
    \\
    \midrule
    Student t-distribution &
    \(\mathbb R\) &
    \(\mu, \sigma\) &
    \(\mathcal T_\nu(x|\mu, \sigma^2)=\dfrac{\Gamma(\frac{\nu+1}{2})}{\sqrt{\pi\nu}\Gamma(\frac{\nu}{2})}\left(1+\dfrac{x^2}{\nu}\right)^{-\frac{\nu+1}{2}}\) &
    \(\mathbb E[X] = 0\) for \(\nu>0\), \(\mathbb V[X] = \frac{\nu}{\nu - 2}\) for \(\nu > 2 \)
    \\
    \midrule
    Cauchy &
    \(\mathbb R\) &
    \(\mu,\lambda\) &
    \(\mathcal C(x|\mu,\lambda)=\dfrac{1}{\lambda\pi(\frac{x-\mu}{\lambda})^2}\) &
    undefined
    \\
    \midrule
    Laplace &
    \(\mathbb R\) &
    \(\mu, b\) &
    \(\text{Laplace}(x|\mu, b) = \dfrac{1}{2b}\exp\left(-\dfrac{|x-\mu|}{b}\right)\) &
    \(\mathbb E[X] = \mu, \mathbb V[X] = 2b^2\)
    \\
    \midrule
    Gamma &
    \(\mathbb R^+\) &
    \(a, b > 0\) &
    Ga($x|\text{shape} = a, \text{scale} = b$) = \(\dfrac{b^a}{\Gamma(a)}x^{a-1}e^{-xb}\) &
    \(\mathbb{E}[X] = ab, \mathbb{V}[X] = ab^2\)
    \\
    \midrule
    Exponential &
    $\mathbb R^+$ &
    $\lambda$ &
    $\text{Expon}(x|\lambda) = \lambda e^{-\lambda x}$ &
    $\mathbb{E}[X] = \lambda, \mathbb{V}[X] = \lambda^2$
    \\
    \midrule
    Chi-squared &
    $\mathbb R^+$ &
    none &
    $\chi_{\nu}^{2}(x) = \newline \text{Ga}(x|\text{shape} = \nu/2, \text{scale} = 2)$ &
    $\mathbb{E}[X] = \nu, \mathbb{V}[X] = 2\nu $
    \\
    \bottomrule
\end{tabularx}
\end{center}

\chapter{Statistics}
This part is about guessing the parameters of our probability models shown in the previous section.
Most of the estimation will take the form of
$\hat{\param} =
\underset{\param}{\text{argmin}}\ \mathcal L(\param)$
, where the $\param$ is our parameter to be estimated, and the $\mathcal L$ is the \emph{loss function}.


There will be two main approaches, namely \textbf{frequentist statistics} and \textbf{Bayesian statistics}.
Frequentist method treats parameters $\param$ as fixed value, and the data  $\mathcal D$ as unknown and changes.
Conversely, Bayesian method treats parameters $\param$ as unknown, and the data  $\mathcal D$ as fixed and known.


However, note that the model may not take any parameters; such models can be adjusted by using \textbf{nonparametric methods}.

    \section{Bayesian Statistics}
    Bayesian statistics treats parameters $\param$ as unknown, and the data $\mathcal D$ as fixed and known.
    We use Bayes rule to update our uncertainty about parameters when the data is given;
    \[
    p(\param|\mathcal D) =
    \dfrac
    {p(\param) p(\mathcal D|\param)}
    {p(\mathcal D)} =
    \dfrac
    {p(\param) p(\mathcal D|\param)}
    {\int p(\param') p(\mathcal D|\param') d\param'}
    \]
    Here, $p(\param)$ is called the \textbf{prior}, and $p(\data|\param)$ is called the \textbf{likelihood}.
    $p(\param|\data)$ is called the \textbf{posterior}, and $p(\data)$ is called the \textbf{marginal likelihood}.

    After the calculation of posterior distribution, we can try to predict the future observations by calculating the \textbf{posterior predictive distribution}:
    \[
    p(x|\data) = \int p(x|\param)p(\param|\data)d\param
    \]
    We are marginalizing out all the unknown parameters.
    If the calculation is too expensive, we can just plug-in the paramter value we estimated, like MLE or MAP.
    \[
    p(x|\data) \approx p(x|\hat{\param})
    \]
    This is called \textbf{plug-in approximation}.

        \subsection{Selecting the Prior}
        Selecting appropriate prior for parameters are one of the challenging problem of Bayesian method.
        There are several choices;
        \begin{itemize}
            \item \hyperref[sec:ConjugatePrior]{conjugate prior}
            \item \hyperref[sec:UninformativePrior]{uninformative prior}
            \item \hyperref[sec:HierPrior]{hierarchical prior}
            \item \hyperref[sec:EmpPrior]{empirical prior}
        \end{itemize}
    
    \section{Frequentist Statistics}
    Frequentist statistics assumes that there is a true parameter value $\param$ and tries to estimate the parameter value
    by changing the data.
    The change of data will change the quantity which is used to estimate the parameter,
    and we use this change to check which is the best estimation.

    Even if the frequentist method has some drawbacks and even paradoxes, several concepts such as
    cross validation is useful even in the Bayesian method. Here, the basic concepts are shown.

        \subsection{Sampling Distribution}
            In frequentist statistics, uncertainty is represented by the \textbf{sampling distribution}
            of a estimator. Here, estimator can be MLE, MAP estimate or method of moments estimate.
            \begin{defn}{Sampling Distribution}{samplingDist}
                Let $\bs{\Theta}(\data)$ be an estimator of parameter $\param$.
                Then, the sampling distribution of $\bs{\Theta}(\data)$ is the distribution obtained by simply keep changing the data set, $\data$.
                More precisely, let $\tilde{\data}^{(s)}$ be an s-th sample from true data distribution. Then,
                \[
                p(\bs{\Theta}(\tilde{\data})=\param) \approx \frac{1}{S}\sum_{s=1}^{S}\delta(\param-\bs{\Theta}(\tilde{\data}^{(s)}))
                \]
                where $S$ is size of the sample.
            \end{defn}

            \subsubsection{Bootstrap Approximation of Sampling Distribution}
                Since the true distribution of data set is unknown, we can't easily get the sampling distribution of estimator.
                Instead, we can use \textbf{bootstrap} method.

                \begin{defn}{Bootstrap Approximation}{bootstapApprox}
                    Let $\data$ be the given data set, and let $p(x|\param^*)$ be the true distribution of data set.
                    \begin{itemize}
                        \item parametric bootstrap: plug-in the estimator value in proability distribution and create sampling distribution.
                            More precisely, use $\tilde{\data}^{(s)} = \{\bs{x}_n \sim p(x|\bs{\Theta}(\data)) | n = 1, 2, \ldots, N\}$
                            instead of $\{\bs{x}_n \sim p(x|\param^*) | n = 1, 2, \ldots, N\}$.
                        \item nonparametric bootstrap: sample datapoints from the original dataset.
                    \end{itemize}
                \end{defn}

            \subsubsection{Bootstrap vs Bayesian}
                An estimator obtained by bootstrap method and posterior distribution are quite different.
                But, in certain condition, where estimator is MLE and the prior for posterior distribution is not so strong,
                they can be quite similar. Check \hyperref[thm:ANSD]{this theorem}.

                Bootstrap method can actually be slower than posterior sampling,
                since bootstrap method requires some number of sampling, but posterior need only one sampling.

            \subsubsection{Asymptotic Normality of Sampling Distribution}
                How can one ensure that the sampling distribution is really useful?
                i.e., can we really approximate real parameter value with sampling distribution?
                The answer is yes, under certain conditions, when the estimator is MLE. 

                \begin{thm}{Asymptotic Normality of Sampling Distribution}{ANSD}
                Under various technical conditions, we have
                \[
                \sqrt{N}(\hat{\param}-\param^*) \rightarrow \mathcal N(\bs{0}, \textbf{F}(\param^*)^{-1})
                \]
                where $\textbf{F}(\param^*)$ is \textbf{Fisher information matrix}, and $\param^*$ is the true parameter value.
                \end{thm}

                Since the \hyperref[defn:FIM]{Fisher information matrix} equals the expectation of Hessian of log likelihood,
                which is shown \hyperref[thm:EBFEHN]{here},
                and Hessian of likelihood measure the amount of curvature of the surface of log likelihood,
                we can infer that the problem with peaked likelihood will give good estimate of parameter when using sampling distribution.

        \subsection{Fisher Information Matrix}
        The Fisher information matrix plays a key role in frequentist statistics,
        and it is also used in Bayesian statistics (Jeffery's prior), and optimization.

            \begin{defn}{Score Function}{ScoreFunc}
            The score function is defined to be the gradient of the log likelihood:
            \[
            \bs{s}(\param) = \nabla_{\param}\log{p(\bs{x}|\param)}
            \]
            \end{defn}
            
            \begin{defn}{Fisher Information Matrix}{FIM}
            Let $\bs{s}(\param)$ be the score function. 
            The Fisher information matrix is defined to be the covariance of the score function:
            \[
            \textbf{F}(\param)=\text{Cov}[\bs{s}(\param) \bs{s}(\param)^{\top}]
            =\mathbb{E}_{\bs{x} \sim p(\bs{x}|\param)}[\nabla_{\param}\log{p(\bs{x}|\param)}\nabla_{\param}\log{p(\bs{x}|\param)}^{\top}]
            \]
            \end{defn}

            And here comes the theorem.

            \begin{thm}{Equivalence Between FIM and Expectation of Hessian of NLL}{EBFEHN}
            Let $\log{p(\bs{x}|\param)}$ be twice-differentiable, and ceratin conditions are satisfied.
            Then, the Fisher information matrix equals the expectation of Hessian of the negative log likelihood,
            i.e.,
            \[
            \mathbb{E}_{\data \sim p(\data|\param)}[\textbf{H}(\data)] = N\textbf{F}(\param)
            \]
            where $\textbf{H}$ is Hessian of $N$ iid samples, $\data$, and $\textbf{F}(\param)$ is the Fisher information matrix.
            \end{thm}

            \begin{proof}
                Search for internet.
            \end{proof}

    \section{Conjugate Priors}\label{sec:ConjugatePrior}
    Calculating the posterior distribution can be untractable, if the prior and likelihood does not match well.
    Here the well-matching prior according to the likelihood is given, called the \textbf{conjugate priors}.

    \begin{center}
    \begin{tabularx}{\textwidth}{XXXXX}
        \toprule
        \textbf{likelihood} &
        \textbf{unknown parameters} &
        \textbf{conjugate prior} &
        \textbf{parameters(prior)} &
        \textbf{posterior}
        \\
        \midrule
            
        \bottomrule
    \end{tabularx}
    \end{center}

    \section{Uninformative Priors}\label{sec:UninformativePrior}
    \section{Hierarchical Priors}\label{sec:HierPrior}
    \section{Empirical Priors}\label{sec:EmpPrior}

\chapter{Decision Theory}
After
\chapter{Graphical Models}
\chapter{Information Theory}
\chapter{Optimization}

\begin{minted}[bgcolor=black!5, fontsize=\small, linenos, frame=lines]{python}
# sample code block
def is_prime(n):
    if n < 2:
        return false
    for i in range(n**0.5, 2):
        if n % i == 0
            return false
        
    return true
\end{minted}

\end{document}